%=========================
% Plantilla LaTeX Artículo Académico
% Compilar con: pdflatex + biber + pdflatex + pdflatex
%=========================
\documentclass[11pt,a4paper]{article}

%-------------------------
% Idioma y tipografía
%-------------------------
\usepackage[english]{babel}
\usepackage{csquotes}
\usepackage[T1]{fontenc}
\usepackage{lmodern}

%-------------------------
% Márgenes y espaciado
%-------------------------
\usepackage[a4paper,margin=2.5cm]{geometry}
\usepackage{setspace}
\onehalfspacing

%-------------------------
% Matemáticas y símbolos
%-------------------------
\usepackage{amsmath,amssymb,amsthm}

%-------------------------
% Tablas y figuras
%-------------------------
\usepackage{graphicx}
\usepackage{booktabs}
\usepackage{caption}
\usepackage{subcaption}
\usepackage{changepage}

%-------------------------
% Enlaces y PDF
%-------------------------
\usepackage{hyperref}
\hypersetup{
  colorlinks=true,
  linkcolor=blue,
  citecolor=blue,
  urlcolor=blue,
  pdftitle={Feferman revisited: conceptual structuralism, Koellner's critique and mathematical practice},
  pdfauthor={Rodrigo Blanco Betes}
}

%-------------------------
% Listas, código y utilidades
%-------------------------
\usepackage{enumitem}
\usepackage{xcolor}
\usepackage{csquotes} % recomendado con biblatex

%-------------------------
% Bibliografía (biblatex + biber)
%-------------------------
\usepackage[
  backend=biber,
  style=apa,
  sortcites=true,
  url=true,
  doi=true
]{biblatex}
\DeclareLanguageMapping{english}{english-apa}
\addbibresource{refs.bib}

%-------------------------
% Metadatos del artículo
%-------------------------
\title{\textbf{Feferman revisited: conceptual structuralism, Koellner's critique and mathematical practice}}

\author{Rodrigo Blanco Betes}


\date{\today}

%-------------------------
% Entornos de teoremas (opcional)
%-------------------------
\newtheorem{theorem}{Theorem}
\newtheorem{lemma}{Lemma}

\begin{document}
\maketitle

\begin{abstract}
Solomon Feferman \textcite{Feferman2011a} argued that CH might not be a well-founded or definite mathematical question. This position has recently been criticized by Koellner \citeyear{Koellner2022} as unfounded, obscure, and overly conventionalist. In this paper, I revisit Feferman’s perspective by paying special attention to mathematical practice. While Feferman’s most discussed argument is related to the clarity of mathematical conceptions, I will emphasize another approach: the critical importance of a mathematical entity being useful for scientific purposes. Interpreted this way, Feferman’s position withstands some of Koellner’s strongest objections but collapses, in my view, when evaluated against the standards of actual mathematical practice. This reveals the need for a more systematic criterion of mathematical legitimacy, one that is grounded in how mathematical entities are used in practice.
\end{abstract}

\section{Introduction}

At the end of the 19th century, Georg Cantor extended the concept of cardinality to infinite sets. In doing so, he was able to show that the cardinality of the continuum is strictly greater than the one of the set of all natural numbers. This led him to conjecture that there was no third set strictly between those two, in terms of cardinality; a conjecture known as the Continuum Hypothesis (CH). The same process can be extended to any set $A$, which can be shown to have smaller cardinality than its power set $\mathcal{P}(A)$. One can then formulate what is known as the Generalized Continuum Hypothesis (GCH), which states that there cannot be an intermediate set strictly between these two in terms of cardinality.\footnote{This intermediate set, in case of existing, would be called an interpolant. There is a stronger way of formulating this hypothesis in terms of well-ordered sets, but it presumes the axiom of choice; see \textcite{Koellner2023}. Even using the weaker version, CH presupposes a huge set theoretical apparatus, consisting in a system strong enough to entail the existence of a definite set containing all the subsets of the continuum, or a strong way of the power set axiom in the case of GCH. In this paper, I focus on the the weaker version, which is also the main focus of Feferman. }

Cantor was unable to resolve this problem, and its status remained open for decades. Years later, \textcite{Godel1940} proved that GCH is consistent with the axioms of ZFC of set theory, and \textcite{Cohen1966} proved the same for its negation. All this gave CH its highly controversial status: an apparently well-defined and legitimate mathematical question but proven to be unsolvable with the tools available to mathematicians in everyday practice.

Since then, there have been a huge number of attempts to settle CH, but none of them has been successful.\footnote{For a summary of the current status of CH, see \textcite{Feferman2011a}, Section 1. } These attempts have never offered a clear counterexample to Cantor's hypothesis, nor an intuitive model demonstrating whether it is true or false. It seems that there is no way to relate CH to a well-established and accepted result of mathematics; so it stands as a metaphysical statement, independent of all the results accepted in everyday mathematical practice.

This has led some mathematicians and philosophers of mathematics to doubt whether it even makes sense to seek a unique and meaningful answer to the continuum hypothesis. One of the most important proponents of this perspective is Solomon Feferman, who proposed a philosophical framework for understanding mathematics called conceptual structuralism. In this framework, mathematics is the result of open-ended processes derived from simple real-life situations, such as counting and ordering, where clarity and practical applicability are crucial. \textcite{Feferman2011a} defended that if we accept conceptual structuralism, then CH turns into an intrinsically vague and indefinite mathematical question. However, conceptual structuralism and its implications for CH have been subject to substantial criticism. Some of the strongest and most direct criticisms have come from Peter Koellner (\citeyear{Koellner2016,Koellner2022}). Koellner's viewpoint on CH bears similarities to those of Woodin (\citeyear{Woodin2001a,Woodin2001b}) or \textcite{Maddy1997}, but he distinguishes himself by having addressed Feferman's philosophical proposal more explicitly.\footnote{\textcite{Parsons2016b} also argues directly against Feferman’s set theoretical conception, but Koellner’s does address Feferman’s case on CH much more explicitly.} Shortly put, Koellner’s main claims are:

\begin{enumerate}
\item That conceptual structuralism is ambiguous on the status of CH.
\item That Feferman’s case relies on a notion of mathematical clarity with several weaknesses, such as the fact that the metamathematical characteristics of set theory that justify Feferman’s position are also present in other areas, such as arithmetic.
\item That it is not acceptable to make mathematics dependent on subjective processes and social conventions as Feferman proposes.
\end{enumerate}

In this paper, I offer an interpretation of conceptual structuralism that makes it easier to understand why it leads to the indefiniteness of CH. I defend that once we consider the emphasis conceptual structuralism places on practical applicability—not much treated in Koellner’s criticism, not even in Feferman’s reasoning—then an argument could be made for the indefiniteness of CH. I do not claim, however, that this argument is definitive, but rather that it is not true that conceptual structuralism is irrelevant to the status of CH. This argument is also, in my view, more promising than Feferman’s idea of clarity. Thus, regarding point (2), I lean more towards Koellner’s view.

A proper treatment of point (3) above would require a deep examination of speech act theory and institutional theory; something that is well beyond the scope of this paper. Let me briefly point out, however, that I do not think, as Koellner does, that it is impossible to understand mathematical entities in an institutional or social way. Quite the contrary. There have been, for the time being, promising attempts to understand mathematics in an institutional way, as Feferman does, and they have proven to be quite resilient to criticism.\footnote{Take for example \textcite{Cole2013}, who gives an institutional account of mathematics based on speech act theory and argues that it can give an account of the necessity and objectivity of mathematical facts.}

In summary, the most important thesis this paper defends is that an alternative interpretation can be given to conceptual structuralism, which focuses most on the applicability of mathematical structures and conceptions in real-life practice, and less on logical and epistemological clarity. I explain in the text, however, that this interpretation, if it honestly confronts its own limitations, does not clearly lead to considering CH as an indefinite or ill-posed mathematical question, as Feferman has often defended.

The text is organized as follows. In the next section, I interpret Feferman’s conceptual structuralism, clarifying and completing it, while also focusing on its implications for CH. Section 3 explains how its relation to practical and scientific applicability could support Feferman’s case, despite presenting limitations. In Section 4, I show how this argument applies to actual mathematical practice. In Sections 5 and 6, I address different ways Feferman approaches the notion of clarity and evaluate the problems his arguments present. Finally, in Section 7 I approach Feferman’s formal argument, and see how it relates to the philosophical ones analyzed in the rest of the text.


\section{Understanding conceptual structuralism}
\label{conceptual-structuralism-intro}

It is not always easy to follow the essential claims of Feferman’s conceptual structuralism. The fundamental idea of this conception is that mathematical notions come from structures based on processes of real experience of activities such as counting, ordering, joining, combining, etc.; and that it is within these structures that mathematical objects and systems are derived and constructed. However, Feferman \citeyear{Feferman2011a} characterizes it in ten theses that are arguably overly concise. He himself acknowledged that they would require further elaboration, but he did not have the space to provide it:

\begin{adjustwidth}{1cm}{1cm}
\begin{quote}
\small
There is no time to elaborate these points, but various aspects of them will come up in our discussion of two constellations of structural notions, first of objects generated by one or more “successor” operations, and second of the continuum. \textcite[p. 12]{Feferman2011a}.
\end{quote}
\end{adjustwidth}

This concerns us here because the most direct criticism by Koellner of Feferman’s argument is precisely that his philosophical conception does not say anything about whether CH must be considered a legitimate or definite statement:

\begin{adjustwidth}{1cm}{1cm}
\begin{quote}
\small
The main point I wish to stress is that conceptual structuralism does not on its own say anything about whether or not a given conception (the natural numbers, arbitrary sets of natural numbers) leads to definite or indefinite statements. We are thus provided with a flexible framework, one that is capable of being applied to views on which CH is definite and also to views on which CH is indefinite. \textcite[p. 10]{Koellner2022}.
\end{quote}
\end{adjustwidth}

For this reason, I would like to give a more complete interpretation of conceptual structuralism: in the first place, by considering Feferman’s other philosophical writings, so that this ambiguity can be addressed. I do not want to literally interpret the grounding theses of conceptual structuralism, but to integrate and understand them within a more complete view of Feferman’s philosophical conception of mathematics. In this regard, I will essentially consider what Feferman says in the whole section 2 of \textcite{Feferman2011a} and in \textcite{Feferman1993}.

In my view, Feferman advances two main lines of argument that bear on the status of CH. First, there is the clarity argument. Here Feferman distinguishes between clear and less clear conceptions of mathematical entities: a clear conception allows statements about it to have determinate truth values, whereas a less clear one can give rise to vague or indeterminate statements, such as the Continuum Hypothesis. This argument integrates naturally with the maxims of conceptual structuralism, since Thesis 5 states:

\begin{adjustwidth}{1cm}{1cm}
\begin{quote}
\small
Basic conceptions differ in their degree of clarity. One may speak of what is true in a given conception, but that notion of truth may only be partial. Truth in full is applicable only to completely clear conceptions. \textcite[p. 11]{Feferman2011a}.
\end{quote}
\end{adjustwidth}

Here, Feferman explicitly acknowledges that truth is not absolute for all conceptions—only fully clear ones can support determinate truth claims. CH, as a statement arising from less clear set-theoretic conceptions, exemplifies this indeterminacy. This argument is the most explicit in Feferman’s case; in fact, Koellner considers it to carry the entire weight of his argumentation:

\begin{adjustwidth}{1cm}{1cm}
\begin{quote}
\small
Finally, in the end, when all the dust settles the entire case rests on the claim that the concept of natural numbers is clear while the concept of arbitrary natural numbers is not clear. \textcite[p. 15]{Koellner2022}.
\end{quote}
\end{adjustwidth}

In \textcite{Feferman2011a}, Section 3, Feferman presents a formal method for deciding whether a mathematical statement is well defined or not, which follows a previous article on semi-constructive theories \textcite{Feferman2010}. He argued that, following this method, CH could perhaps be proven to be an indefinite statement. In these papers Feferman opts for an intermediate approach between classical and intuitionistic logic, using classical logic in some situations he considers suitable for it, and intuitionistic logic in others.

In my view, this logical approach to CH does not constitute a new argument but rather an explanation. Feferman essentially reformulates his philosophical arguments in formal terms, sometimes giving some additional explicitness to its more obscure points. However, I do not think this formal approach gives any further reason to accept Feferman’s case, in case his philosophical arguments are not accepted. Feferman develops the formal system SCS (which stands for Semi-Constructive Set Theory), based on an intuitionistic version of the Kripke–Platek (KP) axiomatic system, with some additional principles added. This results in an axiom system that can be summarized schematically as:

\begin{equation}
  \mathrm{SCS} = \mathrm{IKP} + \mathrm{MP} + \mathrm{AC}
\end{equation}

Here $\mathrm{IKP}$ stands for the intuitionistic Kripke–Platek system with the Axiom of Infinity, $\mathrm{MP}$ for Markov’s principle, and $\mathrm{AC}$ for a constructive, weaker version of the Axiom of Choice.\footnote{See \textcite{Feferman2010}, section 6 for more detail on the technicalities regarding each of the axioms added. The system defined there adds two other principles: the bounded omniscience scheme and the law of excluded middle for $\Delta_0$ formulae, but \textcite{Rathjen2016} has proven them to be redundant.}

One of the most important consequences of these added principles is that the resulting system combines features of both classical and intuitionistic logic. Feferman interprets this to the extent that classical logic can be applied to definite objects and formulas. For example, the validity of the Law of Excluded Middle $A \lor \neg A$ for a given formula $A$ is taken to indicate that $A$ is definite; similarly, other classical principles are tied to definiteness. SCS basically assumes the definiteness—if defined this way—for $\Delta_0$ formulas, which quantify over the existing sets in the theory. In contrast, the universe of sets as a whole is considered indefinite, which is the reason why the system operates in general with intuitionistic logic.

Feferman then came to conjecture that in this axiom system CH was indefinite, that is, that one cannot show $\mathrm{CH} \lor \neg \mathrm{CH}$, so the Law of Excluded Middle is not true for CH in this system. This would entail CH being an indefinite statement within this determinate system. Feferman went further, conjecturing that CH is also indefinite in the system $\mathrm{SCS} + \mathcal{P}(\mathbb{N})$, which adds the existence of the power set of the natural numbers. Both conjectures of Feferman were indeed correct; \textcite{Rathjen2016} published a proof for them. Rathjen, however, was cautious in drawing philosophical conclusions and focused primarily on the technical achievement:

\begin{adjustwidth}{1cm}{1cm}
\begin{quote}
\small
The objective of this paper is to prove Feferman’s conjecture. In this sense it is a technical paper (...). However, as far as the ongoing discussions of the foundational status of CH are concerned, readers will have to form their own conclusions. \textcite[p. 743]{Rathjen2016}.
\end{quote}
\end{adjustwidth}

Here I align with Rathjen: this does not prove Feferman’s philosophical case, but rather clarifies it. Koellner strongly disagrees with this line of argument:

\begin{adjustwidth}{1cm}{1cm}
\begin{quote}
\small
It seems antecedently unsurprising that if one assumes only classical logic for a limited domain—say first-order number theory—then one will not be able to prove that classical logic holds for distinctive statements of a richer domain. \textcite[p. 516]{Koellner2016}.
\end{quote}
\end{adjustwidth}

This reveals the fundamental issue with Feferman’s formal argument if interpreted in this way. Even if CH were shown to be indefinite under certain assumptions, this would not be decisive unless there were independent reasons to accept those assumptions. In this case, Feferman’s assumptions—such as restricting classical reasoning to certain formulas—are justified not by the formal system itself but by his broader philosophical conception of mathematics.

In relation to conceptual structuralism, as discussed earlier, there are two main lines of reasoning that could support the indefiniteness of CH: applicability and clarity. From the perspective of applicability, one might argue that SCS suffices for all scientifically relevant mathematics. From the perspective of clarity, one might claim that SCS captures all mathematical entities that are sufficiently well-defined. In this sense, the formal system does not provide new justification, but rather a more precise articulation of Feferman’s philosophical stance.

There is, however, one aspect of this approach that I find promising: the use of intuitionistic logic to model the notion of indefiniteness. This connects directly with the classical debate between classical and intuitionistic logic. The idea that a formula is definite if and only if it satisfies the Law of Excluded Middle provides a natural bridge between logical principles and philosophical notions of determinacy.

Nevertheless, Feferman’s approach also reveals a tendency toward predicativism, particularly in the exclusion of the full power set operation. In SCS, real numbers are constructed via Cauchy sequences rather than as arbitrary subsets of $\mathbb{N}$. While this suffices for much of analysis, it avoids commitment to the full set-theoretic conception of the continuum.

Feferman justifies this restriction partly on applicability grounds, arguing that all current applications of analysis in physics can be carried out predicatively:

\begin{adjustwidth}{1cm}{1cm}
\begin{quote}
\small
What this shows is that the argument from the success of physics to the reality of the set-theoretical conception of the real numbers is completely vitiated. \textcite[p. 22]{Feferman2011a}.
\end{quote}
\end{adjustwidth}

However, this position seems overly restrictive, as it complicates the development of real analysis and ignores broader aspects of mathematical practice. Feferman himself acknowledges a more moderate position, allowing for the acceptance of $\mathcal{P}(\mathbb{N})$ while rejecting higher levels of the set-theoretic hierarchy. Importantly, his formal results regarding CH remain valid even under this more permissive stance.

%-----------------------------------------

Koellner himself rejects this argument, and, as we will see, there are strong reasons for his stance. For this reason, I think it is reasonable to look for another line of argumentation in Feferman’s thought, that could also lead us to have doubts on the definiteness of CH. This will require both an interpretation and a systematic completion of the philosophical arguments of Feferman, which will result in the applicability argument. This argument defends that mathematical structures are legitimized insofar as they, in some way, model the world and help us understand it. This is not an argument Feferman formulates explicitly, but we see it appear indirectly in Thesis 3:

\begin{adjustwidth}{1cm}{1cm}
\begin{quote}
\small
The basic conceptions of mathematics are of certain kinds of relatively simple ideal world-pictures which are not of objects in isolation but of structures, i.e. coherently conceived groups of objects interconnected by a few simple relations and operations. They are communicated and understood prior to any axiomatics, indeed prior to any systematic logical development. \textcite[p. 11]{Feferman2011a}.
\end{quote}
\end{adjustwidth}

Feferman is likely thinking here of basic conceptions such as the notion of a natural number, the induction principle, and a pre-theoretical idea of the successor operation. From these, some further logical development may be convenient.\footnote{In this regard, it is interesting the idea of unfolding, which refers to taking a concept, definition, or system and systematically developing all of its consequences, rules, or elements in a transparent, step-by-step manner. See \textcite{FefermanStrahm2000}.} These basic conceptions then become part of an ideal world—a coherent conceptual scenario adopted as a framework for doing mathematics. They are “ideal” because they abstract away from physical reality, yet they remain “world pictures” insofar as they provide a useful framework for modeling aspects of the real world. This interpretation of the thesis is grounded in other ideas of Feferman, which I will examine in detail in the following section. It is crucial to note that in formulating this thesis, Feferman talks about “ideal world-pictures”, linking this two last words, while when referring to the conception of the set theoretical hierarchy—see for instance (\citeyear{Feferman2011a}, p. 19)—he refers to an “ideal-world picture”. In (\citeyear{Feferman2011b}, p. 19), he refers to “ideal world pictures”, removing the hyphen, which is probably closer to the first of the other two interpretations. 

Feferman also states here that a mathematical conception consists of a structure, a point that I will adhere to throughout the rest of the text. This means that, whenever I speak of a mathematical conception, I will be referring to sets endowed with corresponding relations, such as the natural numbers with the arithmetical operations or the set-theoretical hierarchy with the belonging relation. It is these structures that will be considered in terms of clarity or applicability, not axiomatic systems.

Taken together, these two arguments could offer complementary perspectives: the applicability argument challenges the legitimacy of frameworks in which CH becomes a question, while the clarity argument challenges the definiteness of CH itself. In line with them, we obtain two corresponding ways of characterizing the status of CH for Feferman. These two lines of thought are distinct but related, since both challenge the idea that CH has a well-defined truth value within a legitimate mathematical conception.

My focus now is to examine this latter argument. I first elaborate on its main ideas and argue that Koellner’s criticism overlooks much of its line of thought. I then turn to the clarity argument. Although the latter is more explicitly discussed in Feferman’s writings, I defend the view that it is the applicability argument that can potentially carry the greatest weight in the case for the indefiniteness of CH, without denying that it also has limitations.

\section{CH regarding the applicability argument}
\label{applicability}

Flow: Feferman implies cognitive limitations to mathematical structures. The question is: how do we \textit{measure} those limitations. One reasonable answer is: mathematical practice and connection to science.

Feferman was in general skeptical about the idea that mathematical entities exist in an abstract, Platonic realm independent of reality. At the same time, he acknowledged that the laws of physics are formulated in a highly idealized way, and that scientific models do not necessarily correspond directly to the real world:

\begin{adjustwidth}{1cm}{1cm}
\begin{quote}
\small
One must ask whether the mathematical structure of the real number system can be identified with the physical structure, or whether it is instead simply an idealized mathematical model of the latter, much as the laws of physics formulated in mathematical terms are highly idealized models of aspects of physical reality. \textcite[p. 107]{Feferman1999}.
\end{quote}
\end{adjustwidth}

Koellner even sees this as a difficulty in Feferman’s perspective, as he notes:

\begin{adjustwidth}{1cm}{1cm}
\begin{quote}
\small
There is a hint of an answer to this question in his remarks on the role of mathematical models in physics. For in addition to thinking that no mythical third realm (the platonic realm) can secure the definiteness of set theory, he also thinks that the physical world cannot secure the definiteness of set theory. \textcite[p. 507]{Koellner2016}.
\end{quote}
\end{adjustwidth}

But the fact that the physical world does not correspond directly to any mathematical conception does not imply that mathematical “idealized” conceptions are completely independent of physical reality. I think what is most coherent with Feferman’s work is to assume that mathematical entities are justified in a pragmatic, not ontological, way; in this regard, the key link lies in our practices—how we model, measure, and reason about the world. This line of reasoning could be considered similar to Quine’s and Putnam’s indispensability argument, but it has important differences from it. Quine and Putnam were more concerned with reasons why the acceptance of the truth of certain physical theories should lead us to commit to corresponding mathematical entities.\footnote{In fact, these arguments are not widely accepted among philosophers of mathematics. See, for instance, a critical formulation and examination in \textcite{Maddy1992}.} For this reason, their focus was epistemological and ontological. By contrast, the applicability argument leans toward a more pragmatic standpoint: consistent use in scientific practice legitimizes a mathematical framework. In general, the relevance of a framework to scientific practice ensures the determinacy of any statement formulated within it; otherwise, it would be difficult to conduct actual science using that framework. For this reason, the fact that scientific practice requires accepting determinate mathematical frameworks legitimizes them.

Feferman conjectures that all mathematics applicable to real-world physics is derivable from simple systems such as Peano Arithmetic:

\begin{adjustwidth}{1cm}{1cm}
\begin{quote}
\small
I have come to conjecture that practically all scientifically applicable mathematics can be formalized in systems reducible to PA, or, as I have sloganized it in \textcite{Feferman1993}: a little bit goes a long way. Against this, I have learned of a couple of cases in some approaches to the foundations of quantum field theory where it appears one must go beyond the resources of PA; but the physical theories that require such additional strength are rather speculative. \textcite[p. 15]{Feferman1999}.
\end{quote}
\end{adjustwidth}

This is no more than a conjecture, but it is a quite reasonable one; most formulations of physics depend either on special relativity, quantum mechanics or general relativity, and neither of these require an especially strong mathematical system. If it is too bold to say that this can be done in PA, a stronger system like $\mathsf{ACA}_0$ is then available. This can be formulated within second-order arithmetic and enables the formation of sets of natural numbers defined by arithmetical formulas. The formalization of calculus, algebra, functional analysis, basic topology, and many mathematical areas that are applied in physics can be carried out through this system. By contrast, statements like the Continuum Hypothesis require systems of third-order arithmetic, which quantify not only over natural numbers and sets of natural numbers (real numbers) but also over sets of sets of natural numbers. This highlights that just stating CH involves a level of set-theoretic complexity far beyond what is needed for the core mathematics of physics, so we have reason to regard this statement as indefinite—at least given the present state of science and scientific practice.

It is worth noting that Koellner’s criticism does not address this line of reasoning. He concentrates on technical and foundational aspects of Feferman’s program and attributes much of its strength to Feferman’s emphasis on clarity. While this is reasonable—since Feferman is much more explicit about clarity—the significance he assigns to practical applicability should not be overlooked in any case.

However, this approach has limitations. As noted, it is reasonable to treat any statement formulated within a mathematical framework fully integrated into scientific practice as definite; otherwise, the framework would be inconvenient for use in science. Arguing the converse—that a statement must belong to such a framework to be mathematically definite—is much more difficult. It is unclear whether Feferman would have argued this way, even when referring to mathematical conceptions as dependent on “ideal world-pictures”. His arguments about applicability were fundamentally negative and were offered in response to the counterargument that non-predicative mathematics is sometimes essential to scientific practice. His positive arguments for the indefiniteness of CH usually depend on the clarity argument, which I will examine in more detail in the following section.

Another problem with this argument, if one were to posit it in favor of the indefiniteness of CH, is its tendency to place too much emphasis on scientific practice, thereby overlooking mathematical practice. Within frameworks that most practitioners consider legitimate, such as the natural numbers, questions may emerge that are not directly relevant to the scientific theories relating to that framework but that are considered interesting by the mathematical community studying them. One such question is Fermat’s Last Theorem. \textcite{Feferman1999} defended the idea that questions like this one would most likely be provable using tools in weaker systems, such as PA. In this particular case, however, the proof required much stronger tools so that, even in the most conservative reconstructions, can only be formulated within sixth-order arithmetic—see \textcite{McLarty2020}. This framework could therefore be legitimized because of its internal applicability to questions that arise within the practice of mathematics. If the idea of a world picture, as formulated in conceptual structuralism, were extended to include this notion of applicability, the conclusion about the indefiniteness of CH could perhaps be completely different.

Finally, I think that Feferman sometimes risks reducing the idea of scientific applicability of mathematical conceptions to their ability to provide the necessary tools for proving relevant mathematical results in physics. Scientific and mathematical practice consists of more than developing a framework in which results can be formulated and proven; they form a network of interplaying practices that involve, for example, practical applications of their results in real life, or trying to disseminate and teach them. In these cases, it is crucial that mathematical frameworks are integrated in the simplest and most natural way into scientific applications. I think that any attempt to impose a condition on mathematical conceptions related to practical applicability must take these ideas into account.

\section{The applicability argument in mathematical practice}

I think that all these limitations must be addressed with a more fine-grained notion of the applicability of mathematical conceptions. In my view, this notion should be consistent with mathematical practice. For many years, philosophy of mathematics has been dominated by the foundationalist tradition, closely linked to foundational programs within mathematics, such as axiomatic set theory. This changed in the latter half of the 20th century with the advent of the so-called philosophy of mathematical practice. As \textcite{Mancosu2008} states:

\begin{adjustwidth}{1cm}{1cm}
\begin{quote}
\small
Mathematical logic, which had been essential in the development of the foundationalist programs, was seen as ineffective in dealing with the questions concerning the dynamics of mathematical discovery and the historical development of mathematics itself. \textcite[p. 4]{Mancosu2008}.
\end{quote}
\end{adjustwidth}

Among many others, some of the most important statements defended by this philosophical tradition are the following:\footnote{We have followed in many ways how Ferreirós addresses the concept of mathematical practice: see (\citeyear{Ferreiros2016}), especially section 2.3.}

\begin{enumerate}
\item Mathematics is a human activity, and not just a compendium of timeless truths. In this activity, reference to the mathematical community is essential.
\item The importance of explanations, clarity and heuristics in mathematics, and not just formal reasoning.
\item Social and historical elements play a crucial role in the development of mathematics. We should not consider that there is necessarily a unique framework that univocally determines the way mathematics is done.
\item Theoretical knowledge is essential to mathematics but must always refer ultimately to practice.
\end{enumerate}

Regardless of whether one accepts the general postulates of this philosophical movement, it is easier to agree that any guiding principle consistently followed in mathematical practice constitutes a central datum that the philosophy of mathematics must account for. Any reasonable completion of conceptual structuralism should try to accommodate such guiding principles.

One of the most influential defenses of methodological principles guiding mathematical practice is the one developed by \textcite{Maddy1997}, who highlights two fundamental complementary maxims: \emph{MAXIMIZE} and \emph{UNIFY}, defending that practitioners tend to make the mathematical universe as large and rich, when consistency is given, as possible. This follows essentially from the principle of quasi-combinatorialism, i.e., the idea that understanding set theory involves grasping the notion of an arbitrary set—definable or not—which ultimately amounts to accepting as many sets as possible. This principle was crucial in the early development of set theory, but it has not continued to be so, either before the advent of mathematical foundationalism or nowadays. What is more, some quasi-combinatorial or maximizing principles are inconsistent with each other; take, for example, the Axiom of Choice and the Axiom of Determinacy. Furthermore, the very concept of an arbitrary set is itself controversial, for instance because it cannot be fully captured within a formal system such as ZFC, or even within much stronger systems, as \textcite{Ferreiros2011} has defended.

Feferman’s conceptual structuralism can in fact be thought of as an alternative to the defense of the “maximize” and “unify” maxims—see \textcite{Feferman2011a}. In this regard, I defend that we should look for similar guiding principles that, in the spirit of the applicability argument, emphasize the role of mathematical entities as world pictures. Perhaps such a guiding principle could be something like the following: all legitimate mathematical conceptions must be essential to formulating results in physics. However, this approach would encounter the limitations discussed at the end of the previous section.

For this reason, we must be more flexible. First, we should accept mathematical conceptions that help us solve internal questions formulated in previously defined concepts or that provide them with explanatory power. Second, we should prefer simpler conceptions, and those that integrate more naturally with the previously defined ones. There is no space here to elaborate on these ideas, which nevertheless naturally extend the guiding principle given before and avoid the issues that arose at the end of the previous section. This amplification of the idea of applicability is what I would call \emph{internal fruitfulness}. I think that an in-depth analysis of this idea, together with its implications for the legitimacy of the mathematical conceptions within which CH can be defined, would be the next step in completing the applicability argument.

Even if this refinement would lead us to reasons for accepting the definiteness of CH, maybe the same could not be said of more abstract mathematical conceptions. Recall that, for CH to be definite, one need only commit to the legitimacy of relatively weak frameworks, described in the language of finite order arithmetic. This lies far behind the full theoretic strength of ZF set theory.

So perhaps the applicability argument, which could be part of a suitable completing interpretation of conceptual structuralism, could be formulated to cast doubt on the legitimacy of higher set-theoretical conceptions, such as the full cumulative hierarchy.

In this regard, the distinction between mathematics and metamathematics is crucial. Axiomatic set theory often functions in this metamathematical role; it provides results about the proving power of specific axiom systems for formalized mathematical languages. A paradigmatic example of a metamathematical reasoning about set-theoretical concepts is found in independence results. These results show the impossibility of logically proving certain statements—most notably, CH—from the axioms of a theory like ZFC. Crucially, these independence proofs do not rely on the abstract set-theoretical structures whose status they investigate. Within the proofs, these structures appear merely as formulas in a formal language. The actual reasoning occurs from an external framework, known as the metatheory. In general, the metatheory is taken to be a much weaker system than ZFC, often requiring only finitistic tools. For example, we see in \textcite{Kunen2011}, p. 11: “All we assert as being ‘really true’ is elementary finitistic reasoning: this is called the metatheory. Using this reasoning, we can explain what a logical formula is and what a formal proof is (…). Then, all discussions about infinite and uncountable sets take place within ZFC, which is called the formal theory”. 

Recall that the standard model of arithmetic can be formalized using the ZF axiom system. Therefore, the consequences of these axioms could not contradict a true result in arithmetic unless the axioms were found to be inconsistent, which seems highly unlikely. For this reason, even a mathematician only committing to the definiteness of arithmetic would find these axioms interesting. They could use the ZF axioms as a useful metamathematical tool and have reasons to believe in the consequences of the axioms when applied to arithmetic. But this would not lead them to commit to the full set-theoretic hierarchy, nor to defend that the concept of an arbitrary set is definite.\footnote{Even acknowledging that ZF is a consistent theory does not imply the aforementioned commitments. A Skolem-type numerable model could suffice for proving the consistency of ZF set theory intuitively.} Whether the set-theoretical hierarchy could be legitimized in the form of “ideal world-pictures,” following the language of Feferman and conceptual structuralism, is not that immediate to address. In this regard, I do not follow Parsons in considering that consistency of axiomatic theories could in general lead to the commitment to a corresponding structure that models them:

\begin{adjustwidth}{1cm}{1cm}
\begin{quote}
\small
We might begin with the weakest assumption that would justify regarding [set-theoretical structures] as legitimate objects of mathematical investigation, the consistency of the basic theories that have been around for some time. \textcite[p. 535]{Parsons2016a}.
\end{quote}
\end{adjustwidth}

In summary, I have exposed the most relevant ideas regarding what I called, in Section 2, the applicability argument. Taking all this into account, I think that the notion of applicability, that Feferman mentions but only indirectly, is nevertheless highly promising. We see here a reasonable criterion for the legitimization of some mathematical conceptions through their appearance in mathematical practice. How CH, or abstract set-theoretical conceptions, might precisely fit into such a framework is a complex topic to think about, that goes beyond the limits of this work (see Author, in preparation).

I think, however, that the other argument, the one of clarity, faces much more difficulties, and could be considered to depend on unstable notions. This is what I would like to address in the following sections.

\section{Considerations around the clarity argument}
\label{clarity}

The clarity argument, as I have framed it, is the line of reasoning Feferman presents more explicitly, both in his philosophical and formal considerations, but that does not prevent it from facing significant difficulties. In summary, Feferman sometimes seems to overstate the lack of clarity in certain mathematical domains. Moreover, the clarity criterion itself can appear unstable, since it depends heavily on individual intuitions. However, these concerns do not undermine the argument entirely.

Feferman’s central idea is that there is a progressive loss of clarity as we move from fully transparent concepts—such as that of a natural number sequence—to increasingly obscure notions:

\begin{adjustwidth}{1cm}{1cm}
\begin{quote}
\small
For example, we have a much clearer conception of arbitrary sequences of points on the Hilbert (or Dedekind, or Cauchy-Cantor) line, or at least of bounded strictly monotone sequence, than we do of arbitrary subsets of the line (...). Our conception allows us to reason that there is no enumeration of all infinite paths through the binary tree, resp. of all subsets of $N$, simply by the appropriate diagonal argument (...). And when we step to $S(S(N))$ there is a still further loss of clarity, but it is just the definiteness of that that is needed to make definite sense of CH. \textcite[p. 21]{Feferman2011a}.
\end{quote}
\end{adjustwidth}

Thesis 5 of conceptual structuralism—see Section 2—maintains that if a conception is not completely clear, then statements formulated within it cannot possess definite truth values. Thus, Feferman’s argument would imply that CH is an indefinite statement, provided his reasons are compelling for regarding higher-order frameworks, such as third-order arithmetic or the cumulative set-theoretic hierarchy (the ones where CH can be precisely formalized), as lacking full clarity.

As just said, this line of argument faces several difficulties. To begin with, Feferman often offers no clear reason for treating certain entities or processes as less well understood than others. For instance, he claims \citeyear{Feferman2011a} that we possess a clearer conception of an infinite path through the full binary tree than of an arbitrary subset of $N$. Here I side with Koellner \citeyear{Koellner2022}, who holds that these notions are essentially equivalent. Feferman also regards the set-theoretic characterization of the continuum as unclear—a point I find highly debatable and will examine in the next two sections. Finally, the idea, given in Thesis 5 of conceptual structuralism, that clarity is a matter of degree raises its own problem: no threshold is given for when diminished clarity makes an entity indeterminate, which suggests that indeterminacy itself would admit degrees, an unintuitive result.

In addition to this, the idea of clarity of a determinate concept seems to be too subjective to constitute a criterion for its validity. This is something Koellner also points out:

\begin{adjustwidth}{1cm}{1cm}
\begin{quote}
\small
Intuitions of clarity are not robust. Some people think that the concept of an arbitrary subset of natural numbers is about as clear a conception as one can have; others disagree. Whose intuitions are to count? \textcite[p. 14]{Koellner2022}.
\end{quote}
\end{adjustwidth}

Feferman himself could be considered to lean towards interpreting clarity as an intrinsically subjective property of a concept, at least that seems to appear in Thesis 6 of conceptual structuralism (\textcite{Feferman2011a}, p. 11): “What is clear in a given conception is time-dependent, both for the individual and historically”. This makes Koellner’s point even stronger.

Feferman’s insistence on clarity can also be interpreted as a modern continuation of the predicative tradition, although he did not always explicitly identify with it:

\begin{adjustwidth}{1cm}{1cm}
\begin{quote}
\small
Though I am not now nor never have been a predicativist, I have to admit to being a sympathizer since I am an avowed anti-platonist, at least insofar as set theory is concerned, and I grant the natural numbers a position of primacy in our mathematical thought. \textcite[p. 313]{Feferman2004a}.
\end{quote}
\end{adjustwidth}

According to predicativism, mathematical entities should be conceptually well-defined to avoid paradoxes or ill-founded reasoning, though the required clarity is now framed in terms of formal, axiomatic structures rather than pre-theoretical intuition alone. Predicativism is hard to reconcile with the idea of defining the continuum set theoretically as a definite totality using the set $\mathcal{P}(\mathbb{N})$. This is essentially because, for a strict predicativist, the power set axiom implies accepting without a clear ground the set of all subsets of a given set, which is a totality that has meaning only in terms of its parts, not necessarily previously defined or built in any satisfactory way.

Let us recall the liar paradox, the statement that “This statement is false”. \textcite{Russell1908} rephrased it as follows: “It is not true for all propositions $p$ that if I affirm $p$, $p$ is true”. He then concluded that the cause of the paradox was a quantification over all the propositions, so that the notion of an arbitrary proposition is illegitimate. A similar reasoning was employed to explain the barber paradox, which appears after trying to define the naïve set ${x \mid x \notin x}$. This concept of an arbitrary naïve set was considered to be unclear, non-definite: this is what was thought to lead to the contradiction. Intuitionists tried to extend this debate even further, by stating that these paradoxes come from accepting logical structures that do not have a counterpart in our intuition (such as the law of excluded middle). We could perhaps say something similar about the concept of an arbitrary subset of the continuum: maybe this concept is so far from our intuition that trying to understand it in linguistic terms may lead us to contradictions or at least incongruities. However, there is no compelling evidence that the concept of an arbitrary subset of the continuum can lead to a contradiction in the same sense as the concept of a naïve set did, or the concept of all propositions. In fact, axiom systems such as ZF are hardly ever doubted to be consistent.

\section{Categoricity and second order logic}
\label{categoricity-second-order-logic}







\footnote{See \textcite{Isaacson2008}, p. 31.} 

If we understand a structure as “clear” only when it is determined up to isomorphism by such a theory, then the universe of sets would indeed lack clarity—but precisely in the same sense in which the natural numbers would lack it. 

We can say the same of the intermediate system of finite order arithmetic.

No recursively axiomatizable first-order theory of arithmetic, including PA, is categorical; Gödel’s incompleteness theorem ensures that categoricity fails at the first-order level. 

One might seek to restore clarity through second-order logic, which does yield categorical axiomatizations of arithmetic and of set-theoretic initial segments. But full second-order semantics presupposes a background notion of “all subsets,” a commitment of the same strength and type as the set-theoretic ontology under discussion. 

Appealing to second-order categoricity to justify the clarity of arithmetic while doubting the clarity of set theory is therefore circular. If we turn our attention to second order logic, we find that $PA_2$ is categorical. 

For set theory, the best result we can have is the quasi-categoricity of $ZC2$ ($ZFC$ without replacement). 

This means that second-order Zermelo set theory with full semantics determines the universe up to an initial segment of the cumulative hierarchy. However, arguing this way that neither set theory of arithmetic are “clear” would be circular, for assuming the same for second order logic. On this point, Feferman and Koellner are in substantial agreement.


\footnote{See p. 19 in \textcite{Feferman2011a} or p. 511 in \textcite{Koellner2016}.}

\textcite{Koellner2010SecondOrderLogic}


The circularity is so patent, one does not see that any more can be claimed thereby than that the definiteness of those notions may be accepted as a matter of faith. \textcite[20]{Feferman2011a}

\subsection{The categoricity of \texorpdfstring{$PA_2$}{PA2}}

The categoricity theorem for second-order Peano Arithmetic ($PA_2$) can be formalized within $ZF\!-\!P+\mathrm{Pow}(\omega)$.\footnote{The axiom $Pow(\omega)$ is understood as the requirement that there exists the power set of the natural numbers.} Under full second-order semantics, second-order variables range over all subsets of the first-order domain. In the case of arithmetic, this amounts to quantification over $\mathcal{P}(\mathbb{N})$. The crucial principle is the second-order induction axiom:

\[
\forall X\Bigl[(0\in X \wedge \forall n\,(n\in X \rightarrow S(n)\in X))
\rightarrow \forall n\,(n\in X)\Bigr].
\]

Once this principle is assumed together with the remaining axioms of $PA_2$, the proof that every model of the theory is isomorphic to the standard model is straightforward.  Be $(\mathbb{N},0,S)$ the standard model of $PA_2$, and suppose that $(M,0_M, S_M)$ is another structure that models this system of axioms. Let $(\mathbb{N},0,S)$ denote the standard model of $PA_2$, and suppose that

\[
(M,0_M,S_M)\models PA_2.
\]

Define a function $f:\mathbb{N}\to M$ recursively by

\[
f(0)=0_M,
\qquad
f(n+1)=S_M(f(n)).
\]

The successor axioms imply that $f$ is injective. To prove surjectivity, define $A=\{x\in M:\exists n\in\mathbb{N}\,(f(n)=x)\}$. By construction, $0_M\in A$. Moreover, if $x\in A$, then for some $n\in\mathbb{N}$ we have $x=f(n)$, and hence $S_M(x) = S_M(f(n)) = f(n+1)$, so $S_M(x)\in A$. Thus $A$ contains $0_M$ and is closed under $S_M$. By the second-order induction axiom, $
A=M$.

Hence every element of $M$ lies in the image of $f$, so $f$ is surjective. Therefore $f$ is a bijection preserving $0$ and successor, and thus an isomorphism. Consequently every full model of $PA_2$ is isomorphic to $(\mathbb{N},0,S)$.

I have highlighted only those aspects of the proof that are relevant to the assumptions on which it depends, rather than its detailed derivation.\footnote{For a rigorous formulation of the proof, see \textcite[82]{shapiro1991foundations}.}

\subsection{The quasi-categoricity of ZFC2}




The crutial point, in my view, is that the set theoretic assumptions made are fairly weak: the most demanding conditions are perhaps accepting the power set of the natural numbers and a fairly weak version of replacement, which tends to happen ubiquitously in mathematical practice:

\begin{adjustwidth}{1cm}{1cm}
\begin{quote}
\small
The categoricity of the $\mathbb{N}$-structure can be obtained in weak second-order logic and does not even depend on its impredicativity. Hence its fully determinate nature, so to speak. In practice, this is admitted by mathematicians and logicians who differ in their acceptance of some questionable foundational principles. \textcite[141]{Ferreiros2023ConStructuralism}
\end{quote}
\end{adjustwidth}

The proof in set theory is more circular than in arithmetic. In either case it is not a mean to stablish definiteness in logical terms.

Definiteness cannot be proven in logical terms, because it trascends logic. It is not to be expected that we can prove logically the definiteness of a certain entity without also departing from the one of a similar one.

\section{Standard models}
\label{standard-models-epistemology}

Non-categoricity implies the existence of non-standard models of both arithmetic and set theory, where certain statements can have different truth values. In arithmetic, every non-standard model contains (as an initial segment of its underlying domain) an image of the standard natural numbers. Therefore, non-standard models are easy to identify. In spite of this, Koellner argues that the case in set theory is exactly the same, and that non-standard models are easy to identify as well:

\begin{adjustwidth}{1cm}{1cm}
\begin{quote}
\small
But the case in set theory is exactly parallel. Take a case like Cohen’s method of forcing where you build a non-standard model (...). This is done within set theory and it is revealed in the first step that one is dealing with a non-standard model since it is a countable object. In the second approach one builds class size models (...) but one builds a Boolean-valued model and, once again, it is immediately revealed that it is non-standard. \textcite[p. 512]{Koellner2016}.
\end{quote}
\end{adjustwidth}

Note how Koellner correctly identifies methods for recognizing a non-standard model of set theory but fails to address how one would positively identify a standard one. This highlights a fundamental disanalogy between arithmetic and set theory. In arithmetic, the standard model, $\mathbb{N}$, is epistemologically privileged. We have a direct pre-theoretical intuition of this structure—it is the very one we use to learn counting. The complete theory of this model serves as a clear (although not recursively computable) benchmark for truth in arithmetic. However, we have no analogous intuition for the entire cumulative set-theoretic hierarchy, where the standard model of set theory is supposed to depend on. Instead, different set-theoretic frameworks propose competing visions of what the universe of sets should be. Most set-theorists would argue that $L$, Gödel’s constructible universe, is not a standard model of the set-theoretic universe. However, to argue this way one might need further philosophical assumptions (such as maximizing principles) which seem much more demanding than arguing that $\mathbb{N}$ is the standard model of arithmetic. The standard model $\mathbb{N}$ is grounded in the simple iterative application of the successor operation, which exhaustively generates the arithmetical universe. By contrast, the cumulative set-theoretic hierarchy is built through a much more intricate network of operations—power set, replacement, transfinite recursion—that do not obviously exhaust the universe of sets.

If we accept conceptual structuralism, a basic conception of mathematics should consist in a structure, from which one could later extract a set of axioms describing it. For that to succeed, there is need of an epistemologically privileged structure, such as $\mathbb{N}$. From this structure one can extract, for example, a set of axioms such as PA, and start reasoning about them. However, the process is inverted in the case of set theory: one begins from a set of axioms, such as ZF, and then tries to accommodate them in a determinate structure, such as the cumulative hierarchy of sets. But as the structure is no longer epistemologically privileged, there is often no way of deciding fundamental questions in it without relying on controversial philosophical principles. I think that, in fact, Feferman is very close to thinking this when he states that “the concept of arbitrary set essential to its formulation is vague or underdetermined and there is no way to sharpen it without violating what it is supposed to be about” (\textcite{Feferman2011a}, p. 1).

This provides a starting point to address the question of how set theory might be considered unclear while arithmetic is not. However, applying this to CH would be problematic because the definiteness of this statement would only correspond to considering third-order arithmetic as clear. It is more difficult to argue that we cannot depart from an intuitive structure, such as $\mathbb{N}$, and develop the corresponding axioms. It is worth recalling that incongruencies such as independence results appear at all levels: arithmetic, finite-order arithmetic, and set theory. \footnote{Koellner treats this in great detail; see (\citeyear{Koellner2016}), sections 5 and 6.}

\textcite{Feferman2000} also considers that set theory can be regarded as having an intuitive model in the cumulative hierarchy, but that does not make him commit to the existence of the entire set-theoretic universe, nor to accepting a Platonic stance:

\begin{adjustwidth}{1cm}{1cm}
\begin{quote}
\small
ZF has an intuitive model in the transfinite iteration of the power set operation taken cumulatively. This has nothing to do with a belief in a platonic reality whose members include the natural numbers and arbitrary sets of natural numbers, and so on. On the contrary, I disbelieve in such entities. \textcite[p. 72]{Feferman2000}.
\end{quote}
\end{adjustwidth}

Perhaps Feferman is distinguishing between a non-standard, yet theoretically valid, model of ZF and a standard model of ZF, which would imply a further commitment to the reality of the universe it describes. He expressed a similar view in a later discussion of CH:

\begin{adjustwidth}{1cm}{1cm}
\begin{quote}
\small
There is no problem to put oneself in the mental frame of mind of “this is what the cumulative hierarchy looks like”, for which one can see that such and such propositions including the axioms of ZFC are (more or less) obviously true. I have taught set theory many times and have presented it in terms of this ideal-world picture with only the caveat that this is what things are supposed to be like in that world, rather than to assert that that’s the way the world actually is. \textcite[p. 19]{Feferman2011a}.
\end{quote}
\end{adjustwidth}

Feferman’s disbelief in a Platonic set-theoretical realm is again evident. His skepticism, however, does not seem to stem from a difficulty in conceiving the realm, but from the distance between such idealized structures and the real world. He concludes the paragraph with a comparison to idealized models in physics:

\begin{adjustwidth}{1cm}{1cm}
\begin{quote}
\small
I am no more uncomfortable doing that than if I were teaching Newtonian mechanics and talking of a world in which space-time is Euclidean (...), or if I were teaching relativistic mechanics. And one can see why little more is needed conceptually to develop an extraordinary part of set theory; as in the case of number theory, a little bit goes a long way. The issue is not whether we can think in such terms but, rather, what the philosophical status is of that kind of thought. \textcite[p. 19]{Feferman2011a}.
\end{quote}
\end{adjustwidth}

In summary, even though Feferman’s requirement for clarity presents significant difficulties, it manages to capture important epistemological differences between arithmetic and set theory. However, extending these differences to third-order arithmetic and the minimal frameworks that allow us to formulate CH would be problematic. This is parallel to the conclusion of Section 4 regarding the applicability argument. In the next section, I will evaluate an important line of thought in Feferman’s case on CH, which I will defend to be essentially a formal reformulation of his philosophical argumentation.

\section{Conclusion}

This is similar to the case around $CH$, where given a countable transitive model $M$ of $ZFC$ we can force to obtain extensions $M[G_0]$, $M[G_0][G_1]$, $M[G_0][G_1][G_2]$,... that alternatively satisfy $CH$ and $\neg CH$. In fact, \textcite{Hamkins2012Multiverse} has used this fact as an argument that $CH$ should be regarded as indefinite.

Be $M \models \mathrm{PA}$ a model of Peano Arithmetic. An end-extension of $M$ is a larger model $N \supseteq M$ such that no element of $N - M$ is inserted ``bellow'' an old element of $M$, i.e., whenever $a \in M$ and $b \in N$ satisfy $b <^N a $, then $b \in M$.

If $\phi$ is an Orey sentence for $\mathrm{PA}$, then if $M$ is a model of $\mathrm{PA}$, we can end-extend it to obtain a model $N_1$ of $\mathrm{PA}$ where $\phi $ holds, and another model $N_2$ one where $\neg \phi$ does. 

The Feferman-Koellner dispute on CH is just one more episode in the long line of philosophical and mathematical controversies that this conjecture has generated since Cantor posed it. Both Feferman and Koellner have made precise and insightful contributions to these controversies. In my view, the only significant weakness in Koellner’s argument was his failure to consider Feferman’s conceptual structuralism more fully. It is true that Feferman’s writings are sometimes too concise and difficult to understand, but I believe the interpretation I have given here is loyal to his main ideas.

The fundamental thesis of this paper is that this philosophical conception, once interpreted as I do, offers a complementary argument for the indefiniteness of CH, which is grounded also in scientific and practical applicability, and not only in the notion of mathematical clarity, upon which Feferman explicitly relies. Considering the possibility of this applicability argument, I believe Koellner is incorrect in stating that conceptual structuralism does not address the indefiniteness of CH. However, this does not mean that we should accept the argument in its basic form, as presented in Section 3.

The next step should be to develop a more fine-grained notion of the applicability of mathematical structures, which includes both the notion of applicability for natural science and the notion of internal applicability to solve problems that arise within previously defined mathematical conceptions.

I argue that this notion can be understood through a recursive sequence, that I call $S_0, S_1, \dots$, which represent ever-widening classes of mathematical conceptions. The class $S_0$ would consist of structures whose legitimacy is grounded directly in their indispensability in modeling the physical world, and would probably include structures such as the set of natural numbers $\mathbb{N}$, the rational numbers $\mathbb{Q}$, the Euclidean plane for spatial reasoning, and others representing core notions of real analysis and calculus (such as that of continuous change). Any structure in $S_0$ can then be justified through the external applicability requirement in its most direct form. Once $S_0$ is defined, it is the moment to take the recursive step. A structure $X$, not belonging to $S = S_1 \cup \dots \cup S_n$, is admitted into class $S_{n+1}$ if it meets, in the most standard, simple and natural way, one of the following conditions:

\begin{enumerate}
  \item \textbf{Problem-Solving Power:} $X$ provides solutions to previously existing, well-defined and relevant problems, i.e., to problems formulated within structures in $S$.

  \item \textbf{Explanatory Unification:} $X$ is able to reveal that previously seemingly disparate phenomena derived from structures in $S$ are in fact manifestations of a single underlying principle from $X$.

  \item \textbf{Generative Capacity:} $X$ naturally leads to new, fruitful lines of inquiry in $S$.
\end{enumerate}

This definition requires further explanation, but I do not have the space to provide it now. The main goal of \textcite{Blanco2} is precisely to further explain what I understand as standard, simple, and natural mathematical conceptions and to state more explicitly which structures could fit into each level of the $\{S_n\}$ hierarchy. In particular, it would shed considerable light on the debate around CH to understand where, within this hierarchy, both systems of third-order arithmetic and more abstract set-theoretic conceptions stand.

My purpose in this paper has been to show how the notion of practical and scientific applicability can play a crucial role in the debate over the definiteness of CH, and how this leads to the next step in understanding such an enigmatic and controversial statement.

%=========================
\printbibliography

\end{document}